
\documentclass{article}
%%%%%%%%%%%%%%%%%%%%%%%%%%%%%%%%%%%%%%%%%%%%%%%%%%%%%%%%%%%%%%%%%%%%%%%%%%%%%%%%%%%%%%%%%%%%%%%%%%%%%%%%%%%%%%%%%%%%%%%%%%%%%%%%%%%%%%%%%%%%%%%%%%%%%%%%%%%%%%%%%%%%%%%%%%%%%%%%%%%%%%%%%%%%%%%%%%%%%%%%%%%%%%%%%%%%%%%%%%%%%%%%%%%%%%%%%%%%%%%%%%%%%%%%%%%%
%TCIDATA{OutputFilter=LATEX.DLL}
%TCIDATA{Version=5.50.0.2960}
%TCIDATA{<META NAME="SaveForMode" CONTENT="1">}
%TCIDATA{BibliographyScheme=Manual}
%TCIDATA{Created=Tuesday, July 14, 2026 15:34:05}
%TCIDATA{LastRevised=Wednesday, July 15, 2026 14:25:16}
%TCIDATA{<META NAME="GraphicsSave" CONTENT="32">}
%TCIDATA{<META NAME="DocumentShell" CONTENT="Scientific Notebook\Blank Document">}
%TCIDATA{CSTFile=Math with theorems suppressed.cst}
%TCIDATA{PageSetup=72,72,72,72,0}
%TCIDATA{AllPages=
%F=36,\PARA{038<p type="texpara" tag="Body Text" >\hfill \thepage}
%}


\newtheorem{theorem}{Theorem}
\newtheorem{acknowledgement}[theorem]{Acknowledgement}
\newtheorem{algorithm}[theorem]{Algorithm}
\newtheorem{axiom}[theorem]{Axiom}
\newtheorem{case}[theorem]{Case}
\newtheorem{claim}[theorem]{Claim}
\newtheorem{conclusion}[theorem]{Conclusion}
\newtheorem{condition}[theorem]{Condition}
\newtheorem{conjecture}[theorem]{Conjecture}
\newtheorem{corollary}[theorem]{Corollary}
\newtheorem{criterion}[theorem]{Criterion}
\newtheorem{definition}[theorem]{Definition}
\newtheorem{example}[theorem]{Example}
\newtheorem{exercise}[theorem]{Exercise}
\newtheorem{lemma}[theorem]{Lemma}
\newtheorem{notation}[theorem]{Notation}
\newtheorem{problem}[theorem]{Problem}
\newtheorem{proposition}[theorem]{Proposition}
\newtheorem{remark}[theorem]{Remark}
\newtheorem{solution}[theorem]{Solution}
\newtheorem{summary}[theorem]{Summary}
\newenvironment{proof}[1][Proof]{\noindent\textbf{#1.} }{\ \rule{0.5em}{0.5em}}
\input{tcilatex}
\begin{document}


\section{An example: The Otto Cycle}

A cyclic process is made from more than one thermodynamic process. The Otto
cycle technically has sic processes. To design an engine we need to describe
all of these processes and make sure they work together to produce a net
useful work.

\FRAME{dtbpF}{4.8888in}{4.0075in}{0pt}{}{}{Figure}{\special{language
"Scientific Word";type "GRAPHIC";maintain-aspect-ratio TRUE;display
"USEDEF";valid_file "T";width 4.8888in;height 4.0075in;depth
0pt;original-width 4.8352in;original-height 3.9583in;cropleft "0";croptop
"1";cropright "1";cropbottom "0";tempfilename
'TI6QJO02.wmf';tempfile-properties "XPR";}}

The $P%
%TCIMACRO{\TeXButton{V}{\ooalign{\hfil$V$\hfil\cr\kern0.1em--\hfil\cr}}}%
%BeginExpansion
\ooalign{\hfil$V$\hfil\cr\kern0.1em--\hfil\cr}%
%EndExpansion
\ $diagram above shows the Otto cycle. This cycle is approximately what
happens in a gasoline engine, like those in cars or, more like what we find
in lawn mowers. It consists of four main processes and two pseudo processes.
All of these processes correspond to a piston moving in a cylinder (our
favorite thermodynamic system!). Here is a picture of what is happening in
the engine. \FRAME{dtbpF}{5.5729in}{1.3589in}{0pt}{}{}{Figure}{\special%
{language "Scientific Word";type "GRAPHIC";maintain-aspect-ratio
TRUE;display "USEDEF";valid_file "T";width 5.5729in;height 1.3589in;depth
0pt;original-width 7.2928in;original-height 1.7553in;cropleft "0";croptop
"1";cropright "1";cropbottom "0";tempfilename
'TI6QJO03.wmf';tempfile-properties "XPR";}}

To understand the Otto cycle let's do the job of finding $\Delta E_{int},$ $%
Q,$ and $w_{int}$ for each part of the cycle and then describe the state of
the gas at the beginning and end of each cycle by finding $P,$ $%
%TCIMACRO{\TeXButton{V}{\ooalign{\hfil$V$\hfil\cr\kern0.1em--\hfil\cr}}}%
%BeginExpansion
\ooalign{\hfil$V$\hfil\cr\kern0.1em--\hfil\cr}%
%EndExpansion
,$ $n,$ $T,$ and $E_{int}$. For our engine, suppose the compression ratio is 
\[
\frac{%
%TCIMACRO{\TeXButton{V}{\ooalign{\hfil$V$\hfil\cr\kern0.1em--\hfil\cr}}}%
%BeginExpansion
\ooalign{\hfil$V$\hfil\cr\kern0.1em--\hfil\cr}%
%EndExpansion
_{A}}{%
%TCIMACRO{\TeXButton{V}{\ooalign{\hfil$V$\hfil\cr\kern0.1em--\hfil\cr}}}%
%BeginExpansion
\ooalign{\hfil$V$\hfil\cr\kern0.1em--\hfil\cr}%
%EndExpansion
_{B}}=8.00 
\]%
Let's take 
\begin{eqnarray*}
%TCIMACRO{\TeXButton{V}{\ooalign{\hfil$V$\hfil\cr\kern0.1em--\hfil\cr}}}%
%BeginExpansion
\ooalign{\hfil$V$\hfil\cr\kern0.1em--\hfil\cr}%
%EndExpansion
_{A} &=&500\unit{cm}^{3}=\allowbreak 0.000\,5\unit{m}^{3} \\
P_{A} &=&100\unit{kPa} \\
T_{A} &=&293\unit{K}
\end{eqnarray*}%
Then we can find $%
%TCIMACRO{\TeXButton{V}{\ooalign{\hfil$V$\hfil\cr\kern0.1em--\hfil\cr}}}%
%BeginExpansion
\ooalign{\hfil$V$\hfil\cr\kern0.1em--\hfil\cr}%
%EndExpansion
_{B}$ 
\[
%TCIMACRO{\TeXButton{V}{\ooalign{\hfil$V$\hfil\cr\kern0.1em--\hfil\cr}}}%
%BeginExpansion
\ooalign{\hfil$V$\hfil\cr\kern0.1em--\hfil\cr}%
%EndExpansion
_{B}=\frac{%
%TCIMACRO{\TeXButton{V}{\ooalign{\hfil$V$\hfil\cr\kern0.1em--\hfil\cr}}}%
%BeginExpansion
\ooalign{\hfil$V$\hfil\cr\kern0.1em--\hfil\cr}%
%EndExpansion
_{A}}{8}=\allowbreak 62.\,\allowbreak 5\unit{cm}^{3}=\allowbreak
6.\,\allowbreak 25\times 10^{-5}\unit{m}^{3} 
\]%
and we will need to measure the temperature at $C.$ Suppose we measure 
\[
T_{C}=1023\unit{K} 
\]%
and the intake temperature to be 
\[
T_{I}=293\unit{K} 
\]%
Further suppose our engine uses a diatomic ideal gas with 
\begin{eqnarray*}
C_{%
%TCIMACRO{\TeXButton{V}{\ooalign{\hfil$V$\hfil\cr\kern0.1em--\hfil\cr}}}%
%BeginExpansion
\ooalign{\hfil$V$\hfil\cr\kern0.1em--\hfil\cr}%
%EndExpansion
} &=&\frac{5}{2}R \\
R &=&8.314\frac{\unit{J}}{\unit{mol}\unit{K}}
\end{eqnarray*}%
so 
\[
\gamma =1.4 
\]%
Let's follow the process through to see what the state variables are at each
point on the PV\ diagram, and what $Q_{h}$ and $Q_{c}$ are for each process,
and what the work, $W,$ is and what the change in internal energy will be, $%
\Delta E_{int}.$ Finally, we will find the efficiency of the engine.

\pagebreak

\subsection{\textbf{Process }$A\rightarrow B$}

\FRAME{dtbpF}{5.0246in}{2.93in}{0pt}{}{}{Figure}{\special{language
"Scientific Word";type "GRAPHIC";maintain-aspect-ratio TRUE;display
"USEDEF";valid_file "T";width 5.0246in;height 2.93in;depth
0pt;original-width 2.9836in;original-height 1.7279in;cropleft "0";croptop
"1";cropright "1";cropbottom "0";tempfilename
'OttoProcessAB_2.wmf';tempfile-properties "XNPR";}}

Assignment:

Fill in the table for process $A\rightarrow B$%
\[
\begin{tabular}{|c|c|c|c|c|c|c|}
\hline
\textbf{State} & $\mathbf{T}\left( \unit{K}\right) $ & $\mathbf{P}\left( 
\unit{kPa}\right) $ & $%
%TCIMACRO{\TeXButton{V}{\ooalign{\hfil$V$\hfil\cr\kern0.1em--\hfil\cr}}}%
%BeginExpansion
\ooalign{\hfil$V$\hfil\cr\kern0.1em--\hfil\cr}%
%EndExpansion
\left( \unit{m}^{3}\right) $ & $\mathbf{n}\left( \unit{mol}\right) $ & $%
\mathbf{E}_{int}\left( \unit{J}\right) $ & \textbf{--} \\ \hline
$\mathbf{A}$ & $293$ & $100$ & $0.000\,5$ & $2.\,052\,5\times 10^{-2}$ &  & 
\textbf{--} \\ \hline
$\mathbf{B}$ &  &  & $6.\,25\times 10^{-5}$ & $2.\,052\,5\times 10^{-2}$ & $%
287.\,\allowbreak 17$ & \textbf{--} \\ \hline
\textbf{Process} & $\mathbf{Q}\left( \unit{J}\right) $ & $\mathbf{w}%
_{int}\left( \unit{J}\right) $ & $\mathbf{\Delta E}_{int}\left( \unit{J}%
\right) $ & $\mathbf{w}_{eng}\left( \unit{J}\right) $ & $\mathbf{Q}_{h}$ & $%
\mathbf{Q}_{c}$ \\ \hline
$\mathbf{AB}$ &  &  &  &  &  &  \\ \hline
\end{tabular}%
\]

ANSWER:

\textbf{Process }$A\rightarrow B$\textbf{\ is an adiabatic compression.} The
piston moves upward, compressing the fuel-air mixture. This happens rapidly,
so there is no time for energy to leave by heat, $Q=0.$ That makes it close
enough to being adiabatic that we can say this process is approximately
adiabatic. In this process $\Delta E_{int}=w_{int}.$

We can use one of our adiabatic equations 
\[
P_{i}%
%TCIMACRO{\TeXButton{V}{\ooalign{\hfil$V$\hfil\cr\kern0.1em--\hfil\cr}}}%
%BeginExpansion
\ooalign{\hfil$V$\hfil\cr\kern0.1em--\hfil\cr}%
%EndExpansion
_{i}^{\gamma }=P_{f}%
%TCIMACRO{\TeXButton{V}{\ooalign{\hfil$V$\hfil\cr\kern0.1em--\hfil\cr}}}%
%BeginExpansion
\ooalign{\hfil$V$\hfil\cr\kern0.1em--\hfil\cr}%
%EndExpansion
_{f}^{\gamma } 
\]%
to find the final pressure for this process at $B$%
\[
P_{B}=P_{A}\frac{%
%TCIMACRO{\TeXButton{V}{\ooalign{\hfil$V$\hfil\cr\kern0.1em--\hfil\cr}}}%
%BeginExpansion
\ooalign{\hfil$V$\hfil\cr\kern0.1em--\hfil\cr}%
%EndExpansion
_{A}^{\gamma }}{%
%TCIMACRO{\TeXButton{V}{\ooalign{\hfil$V$\hfil\cr\kern0.1em--\hfil\cr}}}%
%BeginExpansion
\ooalign{\hfil$V$\hfil\cr\kern0.1em--\hfil\cr}%
%EndExpansion
_{B}^{\gamma }} 
\]%
But of course this works for any point on the line from $A$ to $B$ so 
\[
P=P_{A}\frac{%
%TCIMACRO{\TeXButton{V}{\ooalign{\hfil$V$\hfil\cr\kern0.1em--\hfil\cr}}}%
%BeginExpansion
\ooalign{\hfil$V$\hfil\cr\kern0.1em--\hfil\cr}%
%EndExpansion
_{A}^{\gamma }}{%
%TCIMACRO{\TeXButton{V}{\ooalign{\hfil$V$\hfil\cr\kern0.1em--\hfil\cr}}}%
%BeginExpansion
\ooalign{\hfil$V$\hfil\cr\kern0.1em--\hfil\cr}%
%EndExpansion
^{\gamma }} 
\]

We can use this form for $P$ to find the work for process $A\rightarrow B$ 
\begin{eqnarray*}
w_{AB} &=&-\int_{%
%TCIMACRO{\TeXButton{V}{\ooalign{\hfil$V$\hfil\cr\kern0.1em--\hfil\cr}}}%
%BeginExpansion
\ooalign{\hfil$V$\hfil\cr\kern0.1em--\hfil\cr}%
%EndExpansion
_{A}}^{%
%TCIMACRO{\TeXButton{V}{\ooalign{\hfil$V$\hfil\cr\kern0.1em--\hfil\cr}}}%
%BeginExpansion
\ooalign{\hfil$V$\hfil\cr\kern0.1em--\hfil\cr}%
%EndExpansion
_{B}}Pd%
%TCIMACRO{\TeXButton{V}{\ooalign{\hfil$V$\hfil\cr\kern0.1em--\hfil\cr}} }%
%BeginExpansion
\ooalign{\hfil$V$\hfil\cr\kern0.1em--\hfil\cr}
%EndExpansion
\\
&=&-\int_{%
%TCIMACRO{\TeXButton{V}{\ooalign{\hfil$V$\hfil\cr\kern0.1em--\hfil\cr}}}%
%BeginExpansion
\ooalign{\hfil$V$\hfil\cr\kern0.1em--\hfil\cr}%
%EndExpansion
_{A}}^{%
%TCIMACRO{\TeXButton{V}{\ooalign{\hfil$V$\hfil\cr\kern0.1em--\hfil\cr}}}%
%BeginExpansion
\ooalign{\hfil$V$\hfil\cr\kern0.1em--\hfil\cr}%
%EndExpansion
_{B}}P_{A}\frac{%
%TCIMACRO{\TeXButton{V}{\ooalign{\hfil$V$\hfil\cr\kern0.1em--\hfil\cr}}}%
%BeginExpansion
\ooalign{\hfil$V$\hfil\cr\kern0.1em--\hfil\cr}%
%EndExpansion
_{A}^{\gamma }}{%
%TCIMACRO{\TeXButton{V}{\ooalign{\hfil$V$\hfil\cr\kern0.1em--\hfil\cr}}}%
%BeginExpansion
\ooalign{\hfil$V$\hfil\cr\kern0.1em--\hfil\cr}%
%EndExpansion
^{\gamma }}d%
%TCIMACRO{\TeXButton{V}{\ooalign{\hfil$V$\hfil\cr\kern0.1em--\hfil\cr}} }%
%BeginExpansion
\ooalign{\hfil$V$\hfil\cr\kern0.1em--\hfil\cr}
%EndExpansion
\\
&=&-P_{A}%
%TCIMACRO{\TeXButton{V}{\ooalign{\hfil$V$\hfil\cr\kern0.1em--\hfil\cr}}}%
%BeginExpansion
\ooalign{\hfil$V$\hfil\cr\kern0.1em--\hfil\cr}%
%EndExpansion
_{A}^{\gamma }\int_{%
%TCIMACRO{\TeXButton{V}{\ooalign{\hfil$V$\hfil\cr\kern0.1em--\hfil\cr}}}%
%BeginExpansion
\ooalign{\hfil$V$\hfil\cr\kern0.1em--\hfil\cr}%
%EndExpansion
_{A}}^{%
%TCIMACRO{\TeXButton{V}{\ooalign{\hfil$V$\hfil\cr\kern0.1em--\hfil\cr}}}%
%BeginExpansion
\ooalign{\hfil$V$\hfil\cr\kern0.1em--\hfil\cr}%
%EndExpansion
_{B}}\frac{1}{%
%TCIMACRO{\TeXButton{V}{\ooalign{\hfil$V$\hfil\cr\kern0.1em--\hfil\cr}}}%
%BeginExpansion
\ooalign{\hfil$V$\hfil\cr\kern0.1em--\hfil\cr}%
%EndExpansion
^{\gamma }}d%
%TCIMACRO{\TeXButton{V}{\ooalign{\hfil$V$\hfil\cr\kern0.1em--\hfil\cr}} }%
%BeginExpansion
\ooalign{\hfil$V$\hfil\cr\kern0.1em--\hfil\cr}
%EndExpansion
\\
&=&-P_{A}^{\gamma }%
%TCIMACRO{\TeXButton{V}{\ooalign{\hfil$V$\hfil\cr\kern0.1em--\hfil\cr}}}%
%BeginExpansion
\ooalign{\hfil$V$\hfil\cr\kern0.1em--\hfil\cr}%
%EndExpansion
_{A}^{\gamma }\left( \frac{%
%TCIMACRO{\TeXButton{V}{\ooalign{\hfil$V$\hfil\cr\kern0.1em--\hfil\cr}}}%
%BeginExpansion
\ooalign{\hfil$V$\hfil\cr\kern0.1em--\hfil\cr}%
%EndExpansion
^{1-\gamma }}{1-\gamma }\right\vert _{%
%TCIMACRO{\TeXButton{V}{\ooalign{\hfil$V$\hfil\cr\kern0.1em--\hfil\cr}}}%
%BeginExpansion
\ooalign{\hfil$V$\hfil\cr\kern0.1em--\hfil\cr}%
%EndExpansion
_{A}}^{%
%TCIMACRO{\TeXButton{V}{\ooalign{\hfil$V$\hfil\cr\kern0.1em--\hfil\cr}}}%
%BeginExpansion
\ooalign{\hfil$V$\hfil\cr\kern0.1em--\hfil\cr}%
%EndExpansion
_{B}} \\
&=&-\frac{P_{A}%
%TCIMACRO{\TeXButton{V}{\ooalign{\hfil$V$\hfil\cr\kern0.1em--\hfil\cr}}}%
%BeginExpansion
\ooalign{\hfil$V$\hfil\cr\kern0.1em--\hfil\cr}%
%EndExpansion
_{A}^{\gamma }}{1-\gamma }\left( 
%TCIMACRO{\TeXButton{V}{\ooalign{\hfil$V$\hfil\cr\kern0.1em--\hfil\cr}}}%
%BeginExpansion
\ooalign{\hfil$V$\hfil\cr\kern0.1em--\hfil\cr}%
%EndExpansion
_{B}^{1-\gamma }-%
%TCIMACRO{\TeXButton{V}{\ooalign{\hfil$V$\hfil\cr\kern0.1em--\hfil\cr}}}%
%BeginExpansion
\ooalign{\hfil$V$\hfil\cr\kern0.1em--\hfil\cr}%
%EndExpansion
_{A}^{1-\gamma }\right)
\end{eqnarray*}%
We can rearrange this into a more convenient formula%
\begin{eqnarray*}
w_{AB} &=&\frac{P_{A}%
%TCIMACRO{\TeXButton{V}{\ooalign{\hfil$V$\hfil\cr\kern0.1em--\hfil\cr}}}%
%BeginExpansion
\ooalign{\hfil$V$\hfil\cr\kern0.1em--\hfil\cr}%
%EndExpansion
_{A}^{\gamma }}{\gamma -1}\left( \frac{1}{%
%TCIMACRO{\TeXButton{V}{\ooalign{\hfil$V$\hfil\cr\kern0.1em--\hfil\cr}}}%
%BeginExpansion
\ooalign{\hfil$V$\hfil\cr\kern0.1em--\hfil\cr}%
%EndExpansion
_{B}^{\gamma -1}}-\frac{1}{%
%TCIMACRO{\TeXButton{V}{\ooalign{\hfil$V$\hfil\cr\kern0.1em--\hfil\cr}}}%
%BeginExpansion
\ooalign{\hfil$V$\hfil\cr\kern0.1em--\hfil\cr}%
%EndExpansion
_{A}^{\gamma -1}}\right) \\
&=&\frac{1}{\gamma -1}\left( \frac{P_{A}%
%TCIMACRO{\TeXButton{V}{\ooalign{\hfil$V$\hfil\cr\kern0.1em--\hfil\cr}}}%
%BeginExpansion
\ooalign{\hfil$V$\hfil\cr\kern0.1em--\hfil\cr}%
%EndExpansion
_{A}^{\gamma }}{%
%TCIMACRO{\TeXButton{V}{\ooalign{\hfil$V$\hfil\cr\kern0.1em--\hfil\cr}}}%
%BeginExpansion
\ooalign{\hfil$V$\hfil\cr\kern0.1em--\hfil\cr}%
%EndExpansion
_{B}^{\gamma -1}}-\frac{P_{A}%
%TCIMACRO{\TeXButton{V}{\ooalign{\hfil$V$\hfil\cr\kern0.1em--\hfil\cr}}}%
%BeginExpansion
\ooalign{\hfil$V$\hfil\cr\kern0.1em--\hfil\cr}%
%EndExpansion
_{A}^{\gamma }}{%
%TCIMACRO{\TeXButton{V}{\ooalign{\hfil$V$\hfil\cr\kern0.1em--\hfil\cr}}}%
%BeginExpansion
\ooalign{\hfil$V$\hfil\cr\kern0.1em--\hfil\cr}%
%EndExpansion
_{A}^{\gamma -1}}\right) \\
&=&\frac{1}{\gamma -1}\left( \frac{P_{A}%
%TCIMACRO{\TeXButton{V}{\ooalign{\hfil$V$\hfil\cr\kern0.1em--\hfil\cr}}}%
%BeginExpansion
\ooalign{\hfil$V$\hfil\cr\kern0.1em--\hfil\cr}%
%EndExpansion
_{A}^{\gamma }}{%
%TCIMACRO{\TeXButton{V}{\ooalign{\hfil$V$\hfil\cr\kern0.1em--\hfil\cr}}}%
%BeginExpansion
\ooalign{\hfil$V$\hfil\cr\kern0.1em--\hfil\cr}%
%EndExpansion
_{B}^{\gamma }%
%TCIMACRO{\TeXButton{V}{\ooalign{\hfil$V$\hfil\cr\kern0.1em--\hfil\cr}}}%
%BeginExpansion
\ooalign{\hfil$V$\hfil\cr\kern0.1em--\hfil\cr}%
%EndExpansion
_{B}^{-1}}-\frac{P_{A}%
%TCIMACRO{\TeXButton{V}{\ooalign{\hfil$V$\hfil\cr\kern0.1em--\hfil\cr}}}%
%BeginExpansion
\ooalign{\hfil$V$\hfil\cr\kern0.1em--\hfil\cr}%
%EndExpansion
_{A}^{\gamma }}{%
%TCIMACRO{\TeXButton{V}{\ooalign{\hfil$V$\hfil\cr\kern0.1em--\hfil\cr}}}%
%BeginExpansion
\ooalign{\hfil$V$\hfil\cr\kern0.1em--\hfil\cr}%
%EndExpansion
_{A}^{\gamma -1}}\right)
\end{eqnarray*}%
And once again using 
\[
P_{B}%
%TCIMACRO{\TeXButton{V}{\ooalign{\hfil$V$\hfil\cr\kern0.1em--\hfil\cr}}}%
%BeginExpansion
\ooalign{\hfil$V$\hfil\cr\kern0.1em--\hfil\cr}%
%EndExpansion
_{B}^{\gamma }=P_{A}%
%TCIMACRO{\TeXButton{V}{\ooalign{\hfil$V$\hfil\cr\kern0.1em--\hfil\cr}}}%
%BeginExpansion
\ooalign{\hfil$V$\hfil\cr\kern0.1em--\hfil\cr}%
%EndExpansion
_{A}^{\gamma } 
\]%
or 
\[
P_{A}=P_{B}\frac{%
%TCIMACRO{\TeXButton{V}{\ooalign{\hfil$V$\hfil\cr\kern0.1em--\hfil\cr}}}%
%BeginExpansion
\ooalign{\hfil$V$\hfil\cr\kern0.1em--\hfil\cr}%
%EndExpansion
_{B}^{\gamma }}{%
%TCIMACRO{\TeXButton{V}{\ooalign{\hfil$V$\hfil\cr\kern0.1em--\hfil\cr}}}%
%BeginExpansion
\ooalign{\hfil$V$\hfil\cr\kern0.1em--\hfil\cr}%
%EndExpansion
_{A}^{\gamma }} 
\]%
We get 
\[
w_{AB}=\frac{1}{\gamma -1}\left( \frac{P_{B}\frac{%
%TCIMACRO{\TeXButton{V}{\ooalign{\hfil$V$\hfil\cr\kern0.1em--\hfil\cr}}}%
%BeginExpansion
\ooalign{\hfil$V$\hfil\cr\kern0.1em--\hfil\cr}%
%EndExpansion
_{B}^{\gamma }}{%
%TCIMACRO{\TeXButton{V}{\ooalign{\hfil$V$\hfil\cr\kern0.1em--\hfil\cr}}}%
%BeginExpansion
\ooalign{\hfil$V$\hfil\cr\kern0.1em--\hfil\cr}%
%EndExpansion
_{A}^{\gamma }}%
%TCIMACRO{\TeXButton{V}{\ooalign{\hfil$V$\hfil\cr\kern0.1em--\hfil\cr}}}%
%BeginExpansion
\ooalign{\hfil$V$\hfil\cr\kern0.1em--\hfil\cr}%
%EndExpansion
_{A}^{\gamma }}{%
%TCIMACRO{\TeXButton{V}{\ooalign{\hfil$V$\hfil\cr\kern0.1em--\hfil\cr}}}%
%BeginExpansion
\ooalign{\hfil$V$\hfil\cr\kern0.1em--\hfil\cr}%
%EndExpansion
_{B}^{\gamma }%
%TCIMACRO{\TeXButton{V}{\ooalign{\hfil$V$\hfil\cr\kern0.1em--\hfil\cr}}}%
%BeginExpansion
\ooalign{\hfil$V$\hfil\cr\kern0.1em--\hfil\cr}%
%EndExpansion
_{B}^{-1}}-\frac{P_{A}%
%TCIMACRO{\TeXButton{V}{\ooalign{\hfil$V$\hfil\cr\kern0.1em--\hfil\cr}}}%
%BeginExpansion
\ooalign{\hfil$V$\hfil\cr\kern0.1em--\hfil\cr}%
%EndExpansion
_{A}^{\gamma }}{%
%TCIMACRO{\TeXButton{V}{\ooalign{\hfil$V$\hfil\cr\kern0.1em--\hfil\cr}}}%
%BeginExpansion
\ooalign{\hfil$V$\hfil\cr\kern0.1em--\hfil\cr}%
%EndExpansion
_{A}^{\gamma -1}}\right) 
\]%
A lot of things cancel. We get%
\begin{eqnarray*}
w_{AB} &=&\frac{1}{\gamma -1}\left( \frac{P_{B}}{%
%TCIMACRO{\TeXButton{V}{\ooalign{\hfil$V$\hfil\cr\kern0.1em--\hfil\cr}}}%
%BeginExpansion
\ooalign{\hfil$V$\hfil\cr\kern0.1em--\hfil\cr}%
%EndExpansion
_{B}^{-1}}-\frac{P_{A}%
%TCIMACRO{\TeXButton{V}{\ooalign{\hfil$V$\hfil\cr\kern0.1em--\hfil\cr}}}%
%BeginExpansion
\ooalign{\hfil$V$\hfil\cr\kern0.1em--\hfil\cr}%
%EndExpansion
_{A}^{\gamma }}{%
%TCIMACRO{\TeXButton{V}{\ooalign{\hfil$V$\hfil\cr\kern0.1em--\hfil\cr}}}%
%BeginExpansion
\ooalign{\hfil$V$\hfil\cr\kern0.1em--\hfil\cr}%
%EndExpansion
_{A}^{\gamma -1}}\right) \\
&=&\frac{1}{\gamma -1}\left( P_{B}%
%TCIMACRO{\TeXButton{V}{\ooalign{\hfil$V$\hfil\cr\kern0.1em--\hfil\cr}}}%
%BeginExpansion
\ooalign{\hfil$V$\hfil\cr\kern0.1em--\hfil\cr}%
%EndExpansion
_{B}-P_{A}%
%TCIMACRO{\TeXButton{V}{\ooalign{\hfil$V$\hfil\cr\kern0.1em--\hfil\cr}}}%
%BeginExpansion
\ooalign{\hfil$V$\hfil\cr\kern0.1em--\hfil\cr}%
%EndExpansion
_{A}\right)
\end{eqnarray*}%
which we can generalize into yet another basic equation for adiabatic
processes!%
\[
w_{if}=\left( \frac{1}{\gamma -1}\right) \left( P_{f}%
%TCIMACRO{\TeXButton{V}{\ooalign{\hfil$V$\hfil\cr\kern0.1em--\hfil\cr}}}%
%BeginExpansion
\ooalign{\hfil$V$\hfil\cr\kern0.1em--\hfil\cr}%
%EndExpansion
_{f}-P_{i}%
%TCIMACRO{\TeXButton{V}{\ooalign{\hfil$V$\hfil\cr\kern0.1em--\hfil\cr}}}%
%BeginExpansion
\ooalign{\hfil$V$\hfil\cr\kern0.1em--\hfil\cr}%
%EndExpansion
_{i}\right) 
\]%
Then for process $A\rightarrow B$ we have%
\[
w_{AB}=\left( \frac{1}{\gamma -1}\right) \left( P_{B}%
%TCIMACRO{\TeXButton{V}{\ooalign{\hfil$V$\hfil\cr\kern0.1em--\hfil\cr}}}%
%BeginExpansion
\ooalign{\hfil$V$\hfil\cr\kern0.1em--\hfil\cr}%
%EndExpansion
_{B}-P_{A}%
%TCIMACRO{\TeXButton{V}{\ooalign{\hfil$V$\hfil\cr\kern0.1em--\hfil\cr}}}%
%BeginExpansion
\ooalign{\hfil$V$\hfil\cr\kern0.1em--\hfil\cr}%
%EndExpansion
_{A}\right) 
\]%
and we can use our equation for $P_{B}$ that we got before, $P_{B}=P_{A}%
\frac{%
%TCIMACRO{\TeXButton{V}{\ooalign{\hfil$V$\hfil\cr\kern0.1em--\hfil\cr}}}%
%BeginExpansion
\ooalign{\hfil$V$\hfil\cr\kern0.1em--\hfil\cr}%
%EndExpansion
_{A}^{\gamma }}{%
%TCIMACRO{\TeXButton{V}{\ooalign{\hfil$V$\hfil\cr\kern0.1em--\hfil\cr}}}%
%BeginExpansion
\ooalign{\hfil$V$\hfil\cr\kern0.1em--\hfil\cr}%
%EndExpansion
_{B}^{\gamma }}$ so that our $w_{AB}$ equation becomes 
\[
w_{AB}=\left( \frac{1}{\gamma -1}\right) \left( P_{A}\frac{%
%TCIMACRO{\TeXButton{V}{\ooalign{\hfil$V$\hfil\cr\kern0.1em--\hfil\cr}}}%
%BeginExpansion
\ooalign{\hfil$V$\hfil\cr\kern0.1em--\hfil\cr}%
%EndExpansion
_{A}^{\gamma }}{%
%TCIMACRO{\TeXButton{V}{\ooalign{\hfil$V$\hfil\cr\kern0.1em--\hfil\cr}}}%
%BeginExpansion
\ooalign{\hfil$V$\hfil\cr\kern0.1em--\hfil\cr}%
%EndExpansion
_{B}^{\gamma }}%
%TCIMACRO{\TeXButton{V}{\ooalign{\hfil$V$\hfil\cr\kern0.1em--\hfil\cr}}}%
%BeginExpansion
\ooalign{\hfil$V$\hfil\cr\kern0.1em--\hfil\cr}%
%EndExpansion
_{B}-P_{A}%
%TCIMACRO{\TeXButton{V}{\ooalign{\hfil$V$\hfil\cr\kern0.1em--\hfil\cr}}}%
%BeginExpansion
\ooalign{\hfil$V$\hfil\cr\kern0.1em--\hfil\cr}%
%EndExpansion
_{A}\right) 
\]%
And from our compression ratio equation 
\[
%TCIMACRO{\TeXButton{V}{\ooalign{\hfil$V$\hfil\cr\kern0.1em--\hfil\cr}}}%
%BeginExpansion
\ooalign{\hfil$V$\hfil\cr\kern0.1em--\hfil\cr}%
%EndExpansion
_{B}=\frac{%
%TCIMACRO{\TeXButton{V}{\ooalign{\hfil$V$\hfil\cr\kern0.1em--\hfil\cr}}}%
%BeginExpansion
\ooalign{\hfil$V$\hfil\cr\kern0.1em--\hfil\cr}%
%EndExpansion
_{A}}{8.00} 
\]%
so $w_{AB}$ becomes 
\begin{eqnarray*}
w_{AB} &=&\left( \frac{1}{\gamma -1}\right) \left( P_{A}\frac{%
%TCIMACRO{\TeXButton{V}{\ooalign{\hfil$V$\hfil\cr\kern0.1em--\hfil\cr}}}%
%BeginExpansion
\ooalign{\hfil$V$\hfil\cr\kern0.1em--\hfil\cr}%
%EndExpansion
_{A}^{\gamma }}{\left( \frac{%
%TCIMACRO{\TeXButton{V}{\ooalign{\hfil$V$\hfil\cr\kern0.1em--\hfil\cr}}}%
%BeginExpansion
\ooalign{\hfil$V$\hfil\cr\kern0.1em--\hfil\cr}%
%EndExpansion
_{A}}{8}\right) ^{\gamma }}\frac{%
%TCIMACRO{\TeXButton{V}{\ooalign{\hfil$V$\hfil\cr\kern0.1em--\hfil\cr}}}%
%BeginExpansion
\ooalign{\hfil$V$\hfil\cr\kern0.1em--\hfil\cr}%
%EndExpansion
_{A}}{8.00}-P_{A}%
%TCIMACRO{\TeXButton{V}{\ooalign{\hfil$V$\hfil\cr\kern0.1em--\hfil\cr}}}%
%BeginExpansion
\ooalign{\hfil$V$\hfil\cr\kern0.1em--\hfil\cr}%
%EndExpansion
_{A}\right) \\
&=&\left( \frac{1}{\gamma -1}\right) \left( 8^{\gamma }P_{A}\left( \frac{%
%TCIMACRO{\TeXButton{V}{\ooalign{\hfil$V$\hfil\cr\kern0.1em--\hfil\cr}}}%
%BeginExpansion
\ooalign{\hfil$V$\hfil\cr\kern0.1em--\hfil\cr}%
%EndExpansion
_{A}}{8}\right) -P_{A}%
%TCIMACRO{\TeXButton{V}{\ooalign{\hfil$V$\hfil\cr\kern0.1em--\hfil\cr}}}%
%BeginExpansion
\ooalign{\hfil$V$\hfil\cr\kern0.1em--\hfil\cr}%
%EndExpansion
_{A}\right) \\
&=&\left( \frac{%
%TCIMACRO{\TeXButton{V}{\ooalign{\hfil$V$\hfil\cr\kern0.1em--\hfil\cr}}}%
%BeginExpansion
\ooalign{\hfil$V$\hfil\cr\kern0.1em--\hfil\cr}%
%EndExpansion
_{A}}{\gamma -1}\right) \left( 8^{\gamma -1}P_{A}-P_{A}\right) \\
&=&\left( \frac{%
%TCIMACRO{\TeXButton{V}{\ooalign{\hfil$V$\hfil\cr\kern0.1em--\hfil\cr}}}%
%BeginExpansion
\ooalign{\hfil$V$\hfil\cr\kern0.1em--\hfil\cr}%
%EndExpansion
_{A}P_{A}}{\gamma -1}\right) \left( 8^{\gamma -1}-1\right)
\end{eqnarray*}%
So the work is 
\[
w_{AB}=\left( \frac{%
%TCIMACRO{\TeXButton{V}{\ooalign{\hfil$V$\hfil\cr\kern0.1em--\hfil\cr}}}%
%BeginExpansion
\ooalign{\hfil$V$\hfil\cr\kern0.1em--\hfil\cr}%
%EndExpansion
_{A}P_{A}}{\gamma -1}\right) \left( 8^{\gamma -1}-1\right) 
\]%
and we know all of these numbers so 
\begin{eqnarray*}
w_{AB} &=&\left( \frac{\left( \allowbreak 0.000\,5\unit{m}^{3}\right) \left(
100\unit{kPa}\right) }{1.4-1}\right) \left( 8^{1.4-1}-1\right) \\
&=&\allowbreak 162.\,\allowbreak 17\unit{J}
\end{eqnarray*}

But this is the work done on the gas, the useful work is the negative of this%
\[
w_{eng_{AB}}=-\allowbreak 162.\,\allowbreak 17\unit{J} 
\]%
Which tells us that it took energy to make the process $A\rightarrow B$
happen. You might need a pull cord to start this engine or have an electric
starter motor. The change in internal energy is 
\[
\Delta E_{AB}=\allowbreak 162.\,\allowbreak 17\unit{J} 
\]%
and we can find 
\[
E_{int}=\frac{5}{2}nRT 
\]%
but we need $n.$ We can get it from what we know at $A$ using the ideal gas
law%
\begin{eqnarray*}
n &=&\frac{P_{A}%
%TCIMACRO{\TeXButton{V}{\ooalign{\hfil$V$\hfil\cr\kern0.1em--\hfil\cr}}}%
%BeginExpansion
\ooalign{\hfil$V$\hfil\cr\kern0.1em--\hfil\cr}%
%EndExpansion
_{A}}{RT_{A}} \\
&=&\frac{\left( 100\unit{kPa}\right) \left( \allowbreak 0.000\,5\unit{m}%
^{3}\right) }{\left( 8.314\frac{\unit{J}}{\unit{mol}\unit{K}}\right) \left(
293\unit{K}\right) } \\
&=&\allowbreak 2.\,\allowbreak 052\,5\times 10^{-2}\unit{mol}
\end{eqnarray*}%
Then%
\begin{eqnarray*}
E_{A} &=&\frac{5}{2}nRT_{A} \\
&=&\frac{5}{2}\left( \allowbreak 2.\,\allowbreak 052\,5\times 10^{-2}\unit{%
mol}\right) \left( 8.314\frac{\unit{J}}{\unit{mol}\unit{K}}\right) \left( 293%
\unit{K}\right) \\
&=&\allowbreak 125.\,\allowbreak 00\unit{J}
\end{eqnarray*}

We used 
\[
P_{B}=P_{A}\frac{%
%TCIMACRO{\TeXButton{V}{\ooalign{\hfil$V$\hfil\cr\kern0.1em--\hfil\cr}}}%
%BeginExpansion
\ooalign{\hfil$V$\hfil\cr\kern0.1em--\hfil\cr}%
%EndExpansion
_{A}^{\gamma }}{%
%TCIMACRO{\TeXButton{V}{\ooalign{\hfil$V$\hfil\cr\kern0.1em--\hfil\cr}}}%
%BeginExpansion
\ooalign{\hfil$V$\hfil\cr\kern0.1em--\hfil\cr}%
%EndExpansion
_{B}^{\gamma }} 
\]%
but we should find a numeric answer for $P_{B}.$ Let's use our compression
ratio equation for $%
%TCIMACRO{\TeXButton{V}{\ooalign{\hfil$V$\hfil\cr\kern0.1em--\hfil\cr}}}%
%BeginExpansion
\ooalign{\hfil$V$\hfil\cr\kern0.1em--\hfil\cr}%
%EndExpansion
_{B}$ again 
\begin{eqnarray*}
P_{B} &=&P_{A}\frac{%
%TCIMACRO{\TeXButton{V}{\ooalign{\hfil$V$\hfil\cr\kern0.1em--\hfil\cr}}}%
%BeginExpansion
\ooalign{\hfil$V$\hfil\cr\kern0.1em--\hfil\cr}%
%EndExpansion
_{A}^{\gamma }}{\left( \frac{%
%TCIMACRO{\TeXButton{V}{\ooalign{\hfil$V$\hfil\cr\kern0.1em--\hfil\cr}}}%
%BeginExpansion
\ooalign{\hfil$V$\hfil\cr\kern0.1em--\hfil\cr}%
%EndExpansion
_{A}}{8}\right) ^{\gamma }} \\
&=&P_{A}\frac{%
%TCIMACRO{\TeXButton{V}{\ooalign{\hfil$V$\hfil\cr\kern0.1em--\hfil\cr}}}%
%BeginExpansion
\ooalign{\hfil$V$\hfil\cr\kern0.1em--\hfil\cr}%
%EndExpansion
_{A}^{\gamma }}{\frac{%
%TCIMACRO{\TeXButton{V}{\ooalign{\hfil$V$\hfil\cr\kern0.1em--\hfil\cr}}}%
%BeginExpansion
\ooalign{\hfil$V$\hfil\cr\kern0.1em--\hfil\cr}%
%EndExpansion
_{A}^{\gamma }}{8^{\gamma }}} \\
&=&8^{\gamma }P_{A} \\
&=&8^{1.4}\left( 100\unit{kPa}\right) \\
&=&\allowbreak 1837.\,\allowbreak 9\unit{kPa}
\end{eqnarray*}%
And finally, from the ideal gas law 
\begin{eqnarray*}
T_{B} &=&\frac{P_{B}%
%TCIMACRO{\TeXButton{V}{\ooalign{\hfil$V$\hfil\cr\kern0.1em--\hfil\cr}}}%
%BeginExpansion
\ooalign{\hfil$V$\hfil\cr\kern0.1em--\hfil\cr}%
%EndExpansion
_{B}}{nR} \\
&=&\frac{\left( 1837.\,\allowbreak 9\unit{kPa}\right) \left( 6.\,25\times
10^{-5}\unit{m}^{3}\right) }{\left( 2.\,052\,5\times 10^{-2}\unit{mol}%
\right) \left( 8.314\frac{\unit{J}}{\unit{mol}\unit{K}}\right) } \\
&=&673.\,\allowbreak 15\unit{K}
\end{eqnarray*}%
At last we can find $E_{B}$%
\begin{eqnarray*}
E_{B} &=&\frac{5}{2}nRT_{B} \\
&=&\frac{5}{2}\left( \allowbreak 2.\,\allowbreak 052\,5\times 10^{-2}\unit{%
mol}\right) \left( 8.314\frac{\unit{J}}{\unit{mol}\unit{K}}\right) \left(
673.\,\allowbreak 15\unit{K}\right) \\
&=&\allowbreak 287.\,\allowbreak 17\unit{J}
\end{eqnarray*}%
what we have learned is summarized in this table

\[
\begin{tabular}{|c|c|c|c|c|c|c|}
\hline
\textbf{State} & $\mathbf{T}\left( \unit{K}\right) $ & $\mathbf{P}\left( 
\unit{kPa}\right) $ & $%
%TCIMACRO{\TeXButton{V}{\ooalign{\hfil$V$\hfil\cr\kern0.1em--\hfil\cr}}}%
%BeginExpansion
\ooalign{\hfil$V$\hfil\cr\kern0.1em--\hfil\cr}%
%EndExpansion
\left( \unit{m}^{3}\right) $ & $\mathbf{n}\left( \unit{mol}\right) $ & $%
\mathbf{E}_{int}\left( \unit{J}\right) $ & \textbf{--} \\ \hline
$\mathbf{A}$ & $293$ & $100$ & $0.000\,5$ & $2.\,052\,5\times 10^{-2}$ & $%
\allowbreak 125.\,\allowbreak 00$ & \textbf{--} \\ \hline
$\mathbf{B}$ & $673.\,\allowbreak 15$ & $\allowbreak 1837.\,\allowbreak 9$ & 
$6.\,25\times 10^{-5}$ & $2.\,052\,5\times 10^{-2}$ & $287.\,\allowbreak 17$
& \textbf{--} \\ \hline
\textbf{Process} & $\mathbf{Q}\left( \unit{J}\right) $ & $\mathbf{w}%
_{int}\left( \unit{J}\right) $ & $\mathbf{\Delta E}_{int}\left( \unit{J}%
\right) $ & $\mathbf{w}_{eng}\left( \unit{J}\right) $ & $\mathbf{Q}_{h}$ & $%
\mathbf{Q}_{c}$ \\ \hline
$\mathbf{AB}$ & $0$ & $\allowbreak 162.\,\allowbreak 17\unit{J}$ & $%
\allowbreak 162.\,\allowbreak 17\unit{J}$ & $-\allowbreak 162.\,\allowbreak
17\unit{J}$ & $\mathbf{0}$ & $\mathbf{0}$ \\ \hline
\end{tabular}%
\]

That was just one process!\pagebreak

\subsection{\textbf{Process }$B\rightarrow C$}

Assignment"

Fill in the table for process $B\rightarrow C$

ANSWER:

\[
\begin{tabular}{|c|c|c|c|c|c|c|}
\hline
\textbf{State} & $\mathbf{T}\left( \unit{K}\right) $ & $\mathbf{P}\left( 
\unit{kPa}\right) $ & $%
%TCIMACRO{\TeXButton{V}{\ooalign{\hfil$V$\hfil\cr\kern0.1em--\hfil\cr}}}%
%BeginExpansion
\ooalign{\hfil$V$\hfil\cr\kern0.1em--\hfil\cr}%
%EndExpansion
\left( \unit{m}^{3}\right) $ & $\mathbf{n}\left( \unit{mol}\right) $ & $%
\mathbf{E}_{int}\left( \unit{J}\right) $ & \textbf{--} \\ \hline
$\mathbf{B}$ & $673.\,\allowbreak 15$ & $\allowbreak 1837.\,\allowbreak 9$ & 
$6.\,25\times 10^{-5}$ & $2.\,052\,5\times 10^{-2}$ & $287.\,\allowbreak 17$
& \textbf{--} \\ \hline
$\mathbf{C}$ & $1023$ &  &  & $2.\,052\,5\times 10^{-2}$ &  & \textbf{--} \\ 
\hline
\textbf{Process} & $\mathbf{Q}\left( \unit{J}\right) $ & $\mathbf{w}%
_{int}\left( \unit{J}\right) $ & $\mathbf{\Delta E}_{int}\left( \unit{J}%
\right) $ & $\mathbf{w}_{eng}\left( \unit{J}\right) $ & $\mathbf{Q}_{h}$ & $%
Q_{c}$ \\ \hline
$\mathbf{BC}$ &  &  &  &  &  &  \\ \hline
\end{tabular}%
\]

\textbf{Process }$B\rightarrow C$\textbf{\ is an isochoric process.} \FRAME{%
dtbpF}{3.3304in}{2.9568in}{0pt}{}{}{Figure}{\special{language "Scientific
Word";type "GRAPHIC";maintain-aspect-ratio TRUE;display "USEDEF";valid_file
"T";width 3.3304in;height 2.9568in;depth 0pt;original-width
3.2846in;original-height 2.9136in;cropleft "0";croptop "1";cropright
"1";cropbottom "0";tempfilename 'TI6QJO04.wmf';tempfile-properties "XPR";}}%
This is where the spark plug ignites the air-fuel mixture created by the
fuel injectors. A great deal of energy is released in a hurry. The
temperature jumps up and the pressure does too, but the volume is the same
because this process happens in a flash--literally!

So this must be an isochoric process. We expect $w_{BC}=0,$ so $\Delta
E_{BC}=Q_{BC}.$ The pressure increases, so we expect $Q_{BC}$ to be
positive. That means this energy is entering our engine. This must be part
of $Q_{h}$! We know the temperature at $C$ and we know the volume $%
%TCIMACRO{\TeXButton{V}{\ooalign{\hfil$V$\hfil\cr\kern0.1em--\hfil\cr}}}%
%BeginExpansion
\ooalign{\hfil$V$\hfil\cr\kern0.1em--\hfil\cr}%
%EndExpansion
_{C}=%
%TCIMACRO{\TeXButton{V}{\ooalign{\hfil$V$\hfil\cr\kern0.1em--\hfil\cr}}}%
%BeginExpansion
\ooalign{\hfil$V$\hfil\cr\kern0.1em--\hfil\cr}%
%EndExpansion
_{B}$ so we can find the final pressure%
\begin{eqnarray*}
P_{C} &=&\frac{nRT_{C}}{%
%TCIMACRO{\TeXButton{V}{\ooalign{\hfil$V$\hfil\cr\kern0.1em--\hfil\cr}}}%
%BeginExpansion
\ooalign{\hfil$V$\hfil\cr\kern0.1em--\hfil\cr}%
%EndExpansion
_{C}} \\
&=&\frac{\left( 2.\,052\,5\times 10^{-2}\right) \left( 8.314\frac{\unit{J}}{%
\unit{mol}\unit{K}}\right) \left( 1023\right) }{6.\,25\times 10^{-5}\unit{m}%
^{3}} \\
&=&2.\,\allowbreak 793\,1\times 10^{6}\unit{Pa} \\
&=&2793.\,\allowbreak 1\unit{kPa}
\end{eqnarray*}%
We can find the internal energy at $C$ now%
\begin{eqnarray*}
E_{C} &=&\frac{5}{2}nRT_{C} \\
&=&\frac{5}{2}\left( \allowbreak 2.\,\allowbreak 052\,5\times 10^{-2}\unit{%
mol}\right) \left( 8.314\frac{\unit{J}}{\unit{mol}\unit{K}}\right) \left(
1023\unit{K}\right) \\
&=&\allowbreak 436.\,\allowbreak 42\unit{J}
\end{eqnarray*}%
We can use 
\begin{eqnarray*}
Q &=&nC_{V}\Delta T \\
&=&n\frac{5}{2}R\left( T_{C}-T_{B}\right)
\end{eqnarray*}%
For the process $BC$ to find the energy transferred by heat%
\begin{eqnarray*}
Q_{BC} &=&\left( 2.\,052\,5\times 10^{-2}\unit{mol}\right) \frac{5}{2}\left(
8.314\frac{\unit{J}}{\unit{mol}\unit{K}}\right) \left( 1023\unit{K}%
-673.\,\allowbreak 15\unit{K}\right) \\
&=&149.\,\allowbreak 25\unit{J}
\end{eqnarray*}%
We already mentioned that $Q_{BC}=\Delta E_{int_{BC}}$ and $w_{BC}=0,$ so we
can complete our table of state variables and process variables for this
process. 
\[
\begin{tabular}{|c|c|c|c|c|c|c|}
\hline
\textbf{State} & $\mathbf{T}\left( \unit{K}\right) $ & $\mathbf{P}\left( 
\unit{kPa}\right) $ & $%
%TCIMACRO{\TeXButton{V}{\ooalign{\hfil$V$\hfil\cr\kern0.1em--\hfil\cr}}}%
%BeginExpansion
\ooalign{\hfil$V$\hfil\cr\kern0.1em--\hfil\cr}%
%EndExpansion
\left( \unit{m}^{3}\right) $ & $\mathbf{n}\left( \unit{mol}\right) $ & $%
\mathbf{E}_{int}\left( \unit{J}\right) $ & \textbf{--} \\ \hline
$\mathbf{B}$ & $673.\,\allowbreak 15$ & $\allowbreak 1837.\,\allowbreak 9$ & 
$6.\,25\times 10^{-5}$ & $2.\,052\,5\times 10^{-2}$ & $287.\,\allowbreak 17$
& \textbf{--} \\ \hline
$\mathbf{C}$ & $1023$ & $2793.\,\allowbreak 1$ & $6.\,25\times 10^{-5}$ & $%
2.\,052\,5\times 10^{-2}$ & $436.\,\allowbreak 42$ & \textbf{--} \\ \hline
\textbf{Process} & $\mathbf{Q}\left( \unit{J}\right) $ & $\mathbf{w}%
_{int}\left( \unit{J}\right) $ & $\mathbf{\Delta E}_{int}\left( \unit{J}%
\right) $ & $\mathbf{w}_{eng}\left( \unit{J}\right) $ & $\mathbf{Q}_{h}$ & $%
Q_{c}$ \\ \hline
$\mathbf{BC}$ & $149.\,\allowbreak 25$ & $0$ & $149.\,\allowbreak 25$ & $%
\mathbf{0}$ & $149.\,\allowbreak 25$ & $0$ \\ \hline
\end{tabular}%
\]%
\pagebreak

\subsection{\textbf{Process }$C\rightarrow D$}

Assignment:

Fill in the table for Process $C\rightarrow D$

\[
\begin{tabular}{|c|c|c|c|c|c|c|}
\hline
\textbf{State} & $\mathbf{T}\left( \unit{K}\right) $ & $\mathbf{P}\left( 
\unit{kPa}\right) $ & $%
%TCIMACRO{\TeXButton{V}{\ooalign{\hfil$V$\hfil\cr\kern0.1em--\hfil\cr}}}%
%BeginExpansion
\ooalign{\hfil$V$\hfil\cr\kern0.1em--\hfil\cr}%
%EndExpansion
\left( \unit{m}^{3}\right) $ & $\mathbf{n}\left( \unit{mol}\right) $ & $%
\mathbf{E}_{int}\left( \unit{J}\right) $ & \textbf{--} \\ \hline
$\mathbf{C}$ & $1023$ & $2793.\,\allowbreak 1$ & $6.\,25\times 10^{-5}$ & $%
2.\,052\,5\times 10^{-2}$ & $436.\,\allowbreak 42$ & \textbf{--} \\ \hline
$\mathbf{D}$ &  &  &  & $2.\,052\,5\times 10^{-2}$ &  & \textbf{--} \\ \hline
\textbf{Process} & $\mathbf{Q}\left( \unit{J}\right) $ & $\mathbf{w}%
_{int}\left( \unit{J}\right) $ & $\mathbf{\Delta E}_{int}\left( \unit{J}%
\right) $ & $\mathbf{w}_{eng}\left( \unit{J}\right) $ & $\mathbf{Q}_{h}$ & $%
\mathbf{Q}_{c}$ \\ \hline
$\mathbf{CD}$ &  &  &  &  &  &  \\ \hline
\end{tabular}%
\]

ANSWER:

\FRAME{dtbpF}{4.4227in}{3.1384in}{0pt}{}{}{Figure}{\special{language
"Scientific Word";type "GRAPHIC";maintain-aspect-ratio TRUE;display
"USEDEF";valid_file "T";width 4.4227in;height 3.1384in;depth
0pt;original-width 2.1862in;original-height 1.5437in;cropleft "0";croptop
"1";cropright "1";cropbottom "0";tempfilename
'TI6QJO05.wmf';tempfile-properties "XPR";}}

\textbf{Process }$C\rightarrow D$\textbf{\ is another adiabatic process},
this time an expansion. This is sometimes called the \textquotedblleft power
stroke\textquotedblright\ because the hot, high pressure gas pushes the
piston down. This push is what makes the car or lawn mower (snow blower) go.
We know that $Q_{CD}=0$ so $\Delta E_{CD}=w_{CD}$

To find the ideal gas state variables at $D$ let's start with 
\[
P_{i}%
%TCIMACRO{\TeXButton{V}{\ooalign{\hfil$V$\hfil\cr\kern0.1em--\hfil\cr}}}%
%BeginExpansion
\ooalign{\hfil$V$\hfil\cr\kern0.1em--\hfil\cr}%
%EndExpansion
_{i}^{\gamma }=P_{f}%
%TCIMACRO{\TeXButton{V}{\ooalign{\hfil$V$\hfil\cr\kern0.1em--\hfil\cr}}}%
%BeginExpansion
\ooalign{\hfil$V$\hfil\cr\kern0.1em--\hfil\cr}%
%EndExpansion
_{f}^{\gamma } 
\]%
again, then 
\[
P_{D}=P_{C}\frac{%
%TCIMACRO{\TeXButton{V}{\ooalign{\hfil$V$\hfil\cr\kern0.1em--\hfil\cr}}}%
%BeginExpansion
\ooalign{\hfil$V$\hfil\cr\kern0.1em--\hfil\cr}%
%EndExpansion
_{C}^{\gamma }}{%
%TCIMACRO{\TeXButton{V}{\ooalign{\hfil$V$\hfil\cr\kern0.1em--\hfil\cr}}}%
%BeginExpansion
\ooalign{\hfil$V$\hfil\cr\kern0.1em--\hfil\cr}%
%EndExpansion
_{D}^{\gamma }} 
\]%
but from our PV\ diagram we see $%
%TCIMACRO{\TeXButton{V}{\ooalign{\hfil$V$\hfil\cr\kern0.1em--\hfil\cr}}}%
%BeginExpansion
\ooalign{\hfil$V$\hfil\cr\kern0.1em--\hfil\cr}%
%EndExpansion
_{C}=%
%TCIMACRO{\TeXButton{V}{\ooalign{\hfil$V$\hfil\cr\kern0.1em--\hfil\cr}}}%
%BeginExpansion
\ooalign{\hfil$V$\hfil\cr\kern0.1em--\hfil\cr}%
%EndExpansion
_{B}$ and $%
%TCIMACRO{\TeXButton{V}{\ooalign{\hfil$V$\hfil\cr\kern0.1em--\hfil\cr}}}%
%BeginExpansion
\ooalign{\hfil$V$\hfil\cr\kern0.1em--\hfil\cr}%
%EndExpansion
_{D}=%
%TCIMACRO{\TeXButton{V}{\ooalign{\hfil$V$\hfil\cr\kern0.1em--\hfil\cr}}}%
%BeginExpansion
\ooalign{\hfil$V$\hfil\cr\kern0.1em--\hfil\cr}%
%EndExpansion
_{A}$. This makes sense because the cylinder and piston system maximum and
minimum volumes don't change. So we can write%
\[
P_{D}=P_{C}\frac{%
%TCIMACRO{\TeXButton{V}{\ooalign{\hfil$V$\hfil\cr\kern0.1em--\hfil\cr}}}%
%BeginExpansion
\ooalign{\hfil$V$\hfil\cr\kern0.1em--\hfil\cr}%
%EndExpansion
_{B}^{\gamma }}{%
%TCIMACRO{\TeXButton{V}{\ooalign{\hfil$V$\hfil\cr\kern0.1em--\hfil\cr}}}%
%BeginExpansion
\ooalign{\hfil$V$\hfil\cr\kern0.1em--\hfil\cr}%
%EndExpansion
_{A}^{\gamma }}=8^{-\gamma }P_{C} 
\]%
and knowing this we can find the temperature at $D$ using the ideal gas law%
\[
T_{D}=\frac{P_{D}%
%TCIMACRO{\TeXButton{V}{\ooalign{\hfil$V$\hfil\cr\kern0.1em--\hfil\cr}}}%
%BeginExpansion
\ooalign{\hfil$V$\hfil\cr\kern0.1em--\hfil\cr}%
%EndExpansion
_{D}}{nR} 
\]%
Let's numerically calculate $P_{D}$ and $T_{D}$ at this point.%
\begin{eqnarray*}
P_{D} &=&8^{-1.4}\left( 2793.\,\allowbreak 1\right) \\
&=&151.\,\allowbreak 97\unit{kPa}
\end{eqnarray*}%
and 
\begin{eqnarray*}
T_{D} &=&\frac{\left( 151.\,\allowbreak 97\unit{kPa}\right) \left( 0.000\,5%
\unit{m}^{3}\right) }{\left( 2.\,052\,5\times 10^{-2}\unit{mol}\right)
\left( 8.314\frac{\unit{J}}{\unit{mol}\unit{K}}\right) } \\
&=&445.\,\allowbreak 28\unit{K}
\end{eqnarray*}%
And while we are at it we can complete our state $D$ by calculating $E_{D}$%
\begin{eqnarray*}
E_{D} &=&\frac{5}{2}nRT_{D} \\
&=&\frac{5}{2}\left( \allowbreak 2.\,\allowbreak 052\,5\times 10^{-2}\unit{%
mol}\right) \left( 8.314\frac{\unit{J}}{\unit{mol}\unit{K}}\right) \left(
445.\,\allowbreak 28\unit{K}\right) \\
&=&189.\,\allowbreak 96\unit{J}
\end{eqnarray*}%
Now for the work done, we again use our new adiabatic equation 
\begin{eqnarray*}
w_{CD} &=&\left( \frac{1}{\gamma -1}\right) \left( P_{D}%
%TCIMACRO{\TeXButton{V}{\ooalign{\hfil$V$\hfil\cr\kern0.1em--\hfil\cr}}}%
%BeginExpansion
\ooalign{\hfil$V$\hfil\cr\kern0.1em--\hfil\cr}%
%EndExpansion
_{D}-P_{C}%
%TCIMACRO{\TeXButton{V}{\ooalign{\hfil$V$\hfil\cr\kern0.1em--\hfil\cr}}}%
%BeginExpansion
\ooalign{\hfil$V$\hfil\cr\kern0.1em--\hfil\cr}%
%EndExpansion
_{C}\right) \\
&=&\left( \frac{1}{\gamma -1}\right) \left( 8^{-\gamma }P_{C}%
%TCIMACRO{\TeXButton{V}{\ooalign{\hfil$V$\hfil\cr\kern0.1em--\hfil\cr}}}%
%BeginExpansion
\ooalign{\hfil$V$\hfil\cr\kern0.1em--\hfil\cr}%
%EndExpansion
_{A}-P_{C}%
%TCIMACRO{\TeXButton{V}{\ooalign{\hfil$V$\hfil\cr\kern0.1em--\hfil\cr}}}%
%BeginExpansion
\ooalign{\hfil$V$\hfil\cr\kern0.1em--\hfil\cr}%
%EndExpansion
_{B}\right) \\
&=&\left( \frac{P_{C}}{\gamma -1}\right) \left( 8^{-\gamma }%
%TCIMACRO{\TeXButton{V}{\ooalign{\hfil$V$\hfil\cr\kern0.1em--\hfil\cr}}}%
%BeginExpansion
\ooalign{\hfil$V$\hfil\cr\kern0.1em--\hfil\cr}%
%EndExpansion
_{A}-%
%TCIMACRO{\TeXButton{V}{\ooalign{\hfil$V$\hfil\cr\kern0.1em--\hfil\cr}}}%
%BeginExpansion
\ooalign{\hfil$V$\hfil\cr\kern0.1em--\hfil\cr}%
%EndExpansion
_{B}\right)
\end{eqnarray*}%
all of which we know, so%
\begin{eqnarray*}
W_{CD} &=&\left( \frac{2793.1\unit{kPa}}{1.4-1}\right) \left( 8^{-1.4}\left(
0.000\,5\unit{m}^{3}\right) -6.\,25\times 10^{-5}\unit{m}^{3}\right) \\
&=&-246.\,\allowbreak 46\unit{J}
\end{eqnarray*}%
which completes our set for process $CD$ 
\[
\begin{tabular}{|c|c|c|c|c|c|c|}
\hline
\textbf{State} & $\mathbf{T}\left( \unit{K}\right) $ & $\mathbf{P}\left( 
\unit{kPa}\right) $ & $%
%TCIMACRO{\TeXButton{V}{\ooalign{\hfil$V$\hfil\cr\kern0.1em--\hfil\cr}}}%
%BeginExpansion
\ooalign{\hfil$V$\hfil\cr\kern0.1em--\hfil\cr}%
%EndExpansion
\left( \unit{m}^{3}\right) $ & $\mathbf{n}\left( \unit{mol}\right) $ & $%
\mathbf{E}_{int}\left( \unit{J}\right) $ & \textbf{--} \\ \hline
$\mathbf{C}$ & $1023$ & $2793.\,\allowbreak 1$ & $6.\,25\times 10^{-5}$ & $%
2.\,052\,5\times 10^{-2}$ & $436.\,\allowbreak 42$ & \textbf{--} \\ \hline
$\mathbf{D}$ & $445.\,\allowbreak 28$ & $151.\,\allowbreak 97$ & $0.000\,5$
& $2.\,052\,5\times 10^{-2}$ & $189.\,\allowbreak 96\unit{J}$ & \textbf{--}
\\ \hline
\textbf{Process} & $\mathbf{Q}\left( \unit{J}\right) $ & $\mathbf{w}%
_{int}\left( \unit{J}\right) $ & $\mathbf{\Delta E}_{int}\left( \unit{J}%
\right) $ & $\mathbf{w}_{eng}\left( \unit{J}\right) $ & $\mathbf{Q}_{h}$ & $%
\mathbf{Q}_{c}$ \\ \hline
$\mathbf{CD}$ & $0$ & $-246.\,\allowbreak 46$ & $-246.\,\allowbreak 46$ & $%
246.\,\allowbreak 46$ & $\mathbf{0}$ & $\mathbf{0}$ \\ \hline
\end{tabular}%
\]%
\pagebreak

\subsection{\textbf{Process }$D\rightarrow A$}

Fill out the following table for process $D\rightarrow A$

\[
\begin{tabular}{|c|c|c|c|c|c|c|}
\hline
\textbf{State} & $\mathbf{T}\left( \unit{K}\right) $ & $\mathbf{P}\left( 
\unit{kPa}\right) $ & $%
%TCIMACRO{\TeXButton{V}{\ooalign{\hfil$V$\hfil\cr\kern0.1em--\hfil\cr}}}%
%BeginExpansion
\ooalign{\hfil$V$\hfil\cr\kern0.1em--\hfil\cr}%
%EndExpansion
\left( \unit{m}^{3}\right) $ & $\mathbf{n}\left( \unit{mol}\right) $ & $%
\mathbf{E}_{int}\left( \unit{J}\right) $ & \textbf{--} \\ \hline
$\mathbf{D}$ & $445.\,\allowbreak 28$ & $151.\,\allowbreak 97$ & $0.000\,5$
& $2.\,052\,5\times 10^{-2}$ & $189.\,\allowbreak 96\unit{J}$ & \textbf{--}
\\ \hline
$\mathbf{A}$ & $293$ & $100$ & $0.000\,5$ & $2.\,052\,5\times 10^{-2}$ & $%
\allowbreak 125.\,\allowbreak 00$ & \textbf{--} \\ \hline
\textbf{Process} & $\mathbf{Q}\left( \unit{J}\right) $ & $\mathbf{w}%
_{int}\left( \unit{J}\right) $ & $\mathbf{\Delta E}_{int}\left( \unit{J}%
\right) $ & $\mathbf{w}_{eng}\left( \unit{J}\right) $ & $\mathbf{Q}_{h}$ & $%
\mathbf{Q}_{c}$ \\ \hline
$\mathbf{DA}$ &  &  &  &  &  &  \\ \hline
\end{tabular}%
\]

\textbf{Process }$D\rightarrow A$\textbf{\ is an isochoric process.} We are
nearing the end of the cycle, In process $D\rightarrow A$ a valve opens
letting the pressure drop quickly. \FRAME{dtbpF}{3.0891in}{2.5927in}{0pt}{}{%
}{Figure}{\special{language "Scientific Word";type
"GRAPHIC";maintain-aspect-ratio TRUE;display "USEDEF";valid_file "T";width
3.0891in;height 2.5927in;depth 0pt;original-width 3.6123in;original-height
3.0268in;cropleft "0";croptop "1";cropright "1";cropbottom "0";tempfilename
'TI6QJO06.wmf';tempfile-properties "XPR";}}Since this happens quickly, the
piston does not have time to change position, so the volume does not change.
Process $DA$ is (almost) isochoric. We already know $P_{A},$ $%
%TCIMACRO{\TeXButton{V}{\ooalign{\hfil$V$\hfil\cr\kern0.1em--\hfil\cr}}}%
%BeginExpansion
\ooalign{\hfil$V$\hfil\cr\kern0.1em--\hfil\cr}%
%EndExpansion
_{A},$ $n,$ and $T_{A}$ because that is where we started and this is a
cyclic process. But we need to find the process work, heat and change in
internal energy. Since the volume does not change, $W_{DA}=0$ and $\Delta
E_{DA}=Q_{DA}$ which we can find using 
\begin{eqnarray*}
Q_{DA} &=&nC_{V}\Delta T \\
&=&n\frac{5}{2}R\left( T_{A}-T_{D}\right)
\end{eqnarray*}%
which gives%
\begin{eqnarray*}
Q_{DA} &=&\left( 2.\,052\,5\times 10^{-2}\unit{mol}\right) \frac{5}{2}\left(
8.314\frac{\unit{J}}{\unit{mol}\unit{K}}\right) \left( 293\unit{K}%
-445.\,\allowbreak 28\unit{K}\right) \\
&=&-64.\,\allowbreak 964\unit{J}
\end{eqnarray*}%
Which tells us that we have lost energy from the engine by heat. So $Q_{DA}$
must contribute to $Q_{c}.$ 
\[
\begin{tabular}{|c|c|c|c|c|c|c|}
\hline
\textbf{State} & $\mathbf{T}\left( \unit{K}\right) $ & $\mathbf{P}\left( 
\unit{kPa}\right) $ & $%
%TCIMACRO{\TeXButton{V}{\ooalign{\hfil$V$\hfil\cr\kern0.1em--\hfil\cr}}}%
%BeginExpansion
\ooalign{\hfil$V$\hfil\cr\kern0.1em--\hfil\cr}%
%EndExpansion
\left( \unit{m}^{3}\right) $ & $\mathbf{n}\left( \unit{mol}\right) $ & $%
\mathbf{E}_{int}\left( \unit{J}\right) $ & \textbf{--} \\ \hline
$\mathbf{D}$ & $445.\,\allowbreak 28$ & $151.\,\allowbreak 97$ & $0.000\,5$
& $2.\,052\,5\times 10^{-2}$ & $189.\,\allowbreak 96\unit{J}$ & \textbf{--}
\\ \hline
$\mathbf{A}$ & $293$ & $100$ & $0.000\,5$ & $2.\,052\,5\times 10^{-2}$ & $%
\allowbreak 125.\,\allowbreak 00$ & \textbf{--} \\ \hline
\textbf{Process} & $\mathbf{Q}\left( \unit{J}\right) $ & $\mathbf{w}%
_{int}\left( \unit{J}\right) $ & $\mathbf{\Delta E}_{int}\left( \unit{J}%
\right) $ & $\mathbf{w}_{eng}\left( \unit{J}\right) $ & $\mathbf{Q}_{h}$ & $%
\mathbf{Q}_{c}$ \\ \hline
$\mathbf{DA}$ & $-64.\,\allowbreak 964$ & $0$ & $-64.\,\allowbreak 964$ & $0$
& $0$ & $64.\,\allowbreak 964$ \\ \hline
\end{tabular}%
\]%
\pagebreak

\subsection{Process $A\rightarrow I$ and process $I\rightarrow A$}

Assignment:

Fill out the following table for process $A\rightarrow I$

\[
\begin{tabular}{|c|c|c|c|c|c|c|}
\hline
\textbf{State} & $\mathbf{T}\left( \unit{K}\right) $ & $\mathbf{P}\left( 
\unit{kPa}\right) $ & $%
%TCIMACRO{\TeXButton{V}{\ooalign{\hfil$V$\hfil\cr\kern0.1em--\hfil\cr}}}%
%BeginExpansion
\ooalign{\hfil$V$\hfil\cr\kern0.1em--\hfil\cr}%
%EndExpansion
\left( \unit{m}^{3}\right) $ & $\mathbf{n}\left( \unit{mol}\right) $ & $%
\mathbf{E}_{int}\left( \unit{J}\right) $ & \textbf{--} \\ \hline
$\mathbf{A}$ & $293$ & $100$ & $0.000\,5$ & $2.\,052\,5\times 10^{-2}$ & $%
\allowbreak 125.\,\allowbreak 00$ & \textbf{--} \\ \hline
$\mathbf{I}$ & $293$ & $100$ & $6.\,25\times 10^{-5}$ &  & $\allowbreak $ & 
\textbf{--} \\ \hline
\textbf{Process} & $\mathbf{Q}\left( \unit{J}\right) $ & $\mathbf{w}%
_{int}\left( \unit{J}\right) $ & $\mathbf{\Delta E}_{int}\left( \unit{J}%
\right) $ & $\mathbf{w}_{eng}\left( \unit{J}\right) $ & $\mathbf{Q}_{h}$ & $%
\mathbf{Q}_{c}$ \\ \hline
$\mathbf{AI}$ &  &  &  &  &  &  \\ \hline
\end{tabular}%
\]

\FRAME{dhF}{3.2569in}{2.8323in}{0pt}{}{}{Figure}{\special{language
"Scientific Word";type "GRAPHIC";maintain-aspect-ratio TRUE;display
"USEDEF";valid_file "T";width 3.2569in;height 2.8323in;depth
0pt;original-width 3.2119in;original-height 2.789in;cropleft "0";croptop
"1";cropright "1";cropbottom "0";tempfilename
'TI6QJO07.wmf';tempfile-properties "XPR";}}

But we are not done. The opening of the valve let the spent fuel begin to
leave as exhaust. But we can help remove the exhaust by pushing the piston
up. This is done at constant pressure, because the valve is open to the
exhaust system air pressure. So this is almost an isobaric process. But it's
not quite isobaric because we will lose gas by venting it. The number of
moles, $n,$ will change. We can find the new number of moles using the ideal
gas law%
\[
P_{I}%
%TCIMACRO{\TeXButton{V}{\ooalign{\hfil$V$\hfil\cr\kern0.1em--\hfil\cr}}}%
%BeginExpansion
\ooalign{\hfil$V$\hfil\cr\kern0.1em--\hfil\cr}%
%EndExpansion
_{I}=nRT 
\]%
since $P_{I}=P_{A}$ and $%
%TCIMACRO{\TeXButton{V}{\ooalign{\hfil$V$\hfil\cr\kern0.1em--\hfil\cr}}}%
%BeginExpansion
\ooalign{\hfil$V$\hfil\cr\kern0.1em--\hfil\cr}%
%EndExpansion
_{I}=%
%TCIMACRO{\TeXButton{V}{\ooalign{\hfil$V$\hfil\cr\kern0.1em--\hfil\cr}}}%
%BeginExpansion
\ooalign{\hfil$V$\hfil\cr\kern0.1em--\hfil\cr}%
%EndExpansion
_{B}$ because the piston minimum volume is the same each time the piston
goes up. Then we an use our compression ratio for $%
%TCIMACRO{\TeXButton{V}{\ooalign{\hfil$V$\hfil\cr\kern0.1em--\hfil\cr}}}%
%BeginExpansion
\ooalign{\hfil$V$\hfil\cr\kern0.1em--\hfil\cr}%
%EndExpansion
_{A}$ and $%
%TCIMACRO{\TeXButton{V}{\ooalign{\hfil$V$\hfil\cr\kern0.1em--\hfil\cr}}}%
%BeginExpansion
\ooalign{\hfil$V$\hfil\cr\kern0.1em--\hfil\cr}%
%EndExpansion
_{I}$%
\[
\frac{%
%TCIMACRO{\TeXButton{V}{\ooalign{\hfil$V$\hfil\cr\kern0.1em--\hfil\cr}}}%
%BeginExpansion
\ooalign{\hfil$V$\hfil\cr\kern0.1em--\hfil\cr}%
%EndExpansion
_{A}}{%
%TCIMACRO{\TeXButton{V}{\ooalign{\hfil$V$\hfil\cr\kern0.1em--\hfil\cr}}}%
%BeginExpansion
\ooalign{\hfil$V$\hfil\cr\kern0.1em--\hfil\cr}%
%EndExpansion
_{I}}=8.00 
\]%
or, solving for $%
%TCIMACRO{\TeXButton{V}{\ooalign{\hfil$V$\hfil\cr\kern0.1em--\hfil\cr}}}%
%BeginExpansion
\ooalign{\hfil$V$\hfil\cr\kern0.1em--\hfil\cr}%
%EndExpansion
_{I}$ 
\[
%TCIMACRO{\TeXButton{V}{\ooalign{\hfil$V$\hfil\cr\kern0.1em--\hfil\cr}}}%
%BeginExpansion
\ooalign{\hfil$V$\hfil\cr\kern0.1em--\hfil\cr}%
%EndExpansion
_{I}=8/%
%TCIMACRO{\TeXButton{V}{\ooalign{\hfil$V$\hfil\cr\kern0.1em--\hfil\cr}}}%
%BeginExpansion
\ooalign{\hfil$V$\hfil\cr\kern0.1em--\hfil\cr}%
%EndExpansion
_{A} 
\]%
which we could calculate, but it is just the same as $%
%TCIMACRO{\TeXButton{V}{\ooalign{\hfil$V$\hfil\cr\kern0.1em--\hfil\cr}}}%
%BeginExpansion
\ooalign{\hfil$V$\hfil\cr\kern0.1em--\hfil\cr}%
%EndExpansion
_{B}$

We want $n_{I},$ so let's use the ideal gas law for points $A$ and $I.$ We
get 
\[
P_{A}%
%TCIMACRO{\TeXButton{V}{\ooalign{\hfil$V$\hfil\cr\kern0.1em--\hfil\cr}}}%
%BeginExpansion
\ooalign{\hfil$V$\hfil\cr\kern0.1em--\hfil\cr}%
%EndExpansion
_{A}=n_{A}RT_{A}
\]%
and 
\[
P_{I}%
%TCIMACRO{\TeXButton{V}{\ooalign{\hfil$V$\hfil\cr\kern0.1em--\hfil\cr}}}%
%BeginExpansion
\ooalign{\hfil$V$\hfil\cr\kern0.1em--\hfil\cr}%
%EndExpansion
_{I}=n_{I}RT_{I}
\]%
We can divide these two equations%
\[
\frac{P_{A}%
%TCIMACRO{\TeXButton{V}{\ooalign{\hfil$V$\hfil\cr\kern0.1em--\hfil\cr}}}%
%BeginExpansion
\ooalign{\hfil$V$\hfil\cr\kern0.1em--\hfil\cr}%
%EndExpansion
_{A}}{P_{I}%
%TCIMACRO{\TeXButton{V}{\ooalign{\hfil$V$\hfil\cr\kern0.1em--\hfil\cr}}}%
%BeginExpansion
\ooalign{\hfil$V$\hfil\cr\kern0.1em--\hfil\cr}%
%EndExpansion
_{I}}=\frac{n_{A}RT_{A}}{n_{I}RT_{I}}
\]%
and note that we have our compression ratio hidden in the left hand side,
and recall that $P_{I}=P_{A}$ so the left hand side is just $8.$ And on the
right hand side $T_{A}=T_{I}$ because we didn't change the temperature of
the gass, we just moved it. 
\[
8=\frac{n_{A}}{n_{I}}
\]%
We can solve this for the number of moles at piston $I$%
\[
n_{I}=\frac{n_{A}}{8}
\]%
And we have all the numbers, so 
\begin{eqnarray*}
n_{I} &=&\frac{\left( 2.\,052\,5\times 10^{-2}\unit{mol}\right) }{8} \\
&=&\allowbreak 2.\,\allowbreak 565\,6\times 10^{-3}\unit{mol}
\end{eqnarray*}%
And then $E_{I}$ will be 
\begin{eqnarray*}
E_{D} &=&\frac{5}{2}n_{I}RT_{I} \\
&=&\frac{5}{2}\left( 2.\,\allowbreak 565\,6\times 10^{-3}\unit{mol}\right)
\left( 8.314\frac{\unit{J}}{\unit{mol}\unit{K}}\right) \left( 293\unit{K}%
\right)  \\
&=&15.\,\allowbreak 625\unit{J}
\end{eqnarray*}%
The work done for a true isobaric process is 
\[
w_{AI}=P_{A}\left( 
%TCIMACRO{\TeXButton{V}{\ooalign{\hfil$V$\hfil\cr\kern0.1em--\hfil\cr}}}%
%BeginExpansion
\ooalign{\hfil$V$\hfil\cr\kern0.1em--\hfil\cr}%
%EndExpansion
_{I}-%
%TCIMACRO{\TeXButton{V}{\ooalign{\hfil$V$\hfil\cr\kern0.1em--\hfil\cr}}}%
%BeginExpansion
\ooalign{\hfil$V$\hfil\cr\kern0.1em--\hfil\cr}%
%EndExpansion
_{A}\right) 
\]%
but this assumed that $n$ did not change. Really we are just moving the gas,
and it doesn't take much of a force to move the little bit of gas. So to a
good approximation we can say 
\begin{eqnarray*}
w_{int} &=&F\Delta x \\
&\approx &0
\end{eqnarray*}%
and we can ignore $w_{AI}.$

For finding $Q_{AI}$ we would like to use 
\[
Q=nC_{P}\Delta T
\]%
but we quickly run into trouble again because $n$ changed! We are using
convection to remove heat instead of conduction. We can get around this by
finding $\Delta E_{int}$ directly. 
\begin{eqnarray*}
\Delta E_{AI} &=&E_{I}-E_{A} \\
&=&15.\,\allowbreak 625\unit{J}-125.\,\allowbreak 00\unit{J} \\
&=&-109.\,\allowbreak 38\unit{J}
\end{eqnarray*}%
From this and the first law%
\[
\Delta E_{int}=Q+w_{int}
\]%
We can find $Q$%
\begin{eqnarray*}
Q &\approx &\Delta E_{int} \\
&\approx &-109.\,\allowbreak 38\unit{J}
\end{eqnarray*}

Energy is leaving with the gas, so we seem to pick up a contribution to $%
Q_{c}.$ But really this is not part of our thermodynamic process, because we
mechanically removed the energy by removing the gas. Let's denote this by
putting our energy in parenthesis to remind us that this is a more
mechanical than thermal process.%
\[
\begin{tabular}{|c|c|c|c|c|c|c|}
\hline
\textbf{State} & $\mathbf{T}\left( \unit{K}\right) $ & $\mathbf{P}\left( 
\unit{kPa}\right) $ & $%
%TCIMACRO{\TeXButton{V}{\ooalign{\hfil$V$\hfil\cr\kern0.1em--\hfil\cr}}}%
%BeginExpansion
\ooalign{\hfil$V$\hfil\cr\kern0.1em--\hfil\cr}%
%EndExpansion
\left( \unit{m}^{3}\right) $ & $\mathbf{n}\left( \unit{mol}\right) $ & $%
\mathbf{E}_{int}\left( \unit{J}\right) $ & \textbf{--} \\ \hline
$\mathbf{A}$ & $293$ & $100$ & $0.000\,5$ & $2.\,052\,5\times 10^{-2}$ & $%
\allowbreak 125.\,\allowbreak 00$ & \textbf{--} \\ \hline
$\mathbf{I}$ & $293$ & $100$ & $6.\,25\times 10^{-5}$ & $2.\,\allowbreak
565\,6\times 10^{-3}$ & $15.\,\allowbreak 625\unit{J}$ & \textbf{--} \\ 
\hline
\textbf{Process} & $\mathbf{Q}\left( \unit{J}\right) $ & $\mathbf{w}%
_{int}\left( \unit{J}\right) $ & $\mathbf{\Delta E}_{int}\left( \unit{J}%
\right) $ & $\mathbf{w}_{eng}\left( \unit{J}\right) $ & $\mathbf{Q}_{h}$ & $%
\mathbf{Q}_{c}$ \\ \hline
$\mathbf{AI}$ & $\left( -109.\,\allowbreak 38\unit{J}\right) $ & $\sim 0$ & $%
\left( -109.\,\allowbreak 38\right) $ & $\sim 0$ & $0$ & $\left(
109.\,\allowbreak 38\right) $ \\ \hline
\end{tabular}%
\]%
$\allowbreak $\FRAME{dhF}{3.0606in}{3.2932in}{0pt}{}{}{Figure}{\special%
{language "Scientific Word";type "GRAPHIC";maintain-aspect-ratio
TRUE;display "USEDEF";valid_file "T";width 3.0606in;height 3.2932in;depth
0pt;original-width 3.0173in;original-height 3.2474in;cropleft "0";croptop
"1";cropright "1";cropbottom "0";tempfilename
'TI8GQ400.wmf';tempfile-properties "XPR";}}Now the exhaust valve closes, and
the intake valve opens. The new fuel mixture is brought in. We essentially
reverse process $AI$ in process $IA.$ 
\[
\begin{tabular}{|c|c|c|c|c|c|c|}
\hline
\textbf{State} & $\mathbf{T}\left( \unit{K}\right) $ & $\mathbf{P}\left( 
\unit{kPa}\right) $ & $%
%TCIMACRO{\TeXButton{V}{\ooalign{\hfil$V$\hfil\cr\kern0.1em--\hfil\cr}}}%
%BeginExpansion
\ooalign{\hfil$V$\hfil\cr\kern0.1em--\hfil\cr}%
%EndExpansion
\left( \unit{m}^{3}\right) $ & $\mathbf{n}\left( \unit{mol}\right) $ & $%
\mathbf{E}_{int}\left( \unit{J}\right) $ & \textbf{--} \\ \hline
$\mathbf{I}$ & $293$ & $100$ & $6.\,25\times 10^{-5}$ & $2.\,\allowbreak
565\,7\times 10^{-3}$ & $15.\,\allowbreak 625\unit{J}$ & \textbf{--} \\ 
\hline
$\mathbf{A}$ & $293$ & $100$ & $0.000\,5$ & $2.\,052\,5\times 10^{-2}$ & $%
\allowbreak 125.\,\allowbreak 00$ & \textbf{--} \\ \hline
\textbf{Process} & $\mathbf{Q}\left( \unit{J}\right) $ & $\mathbf{w}%
_{int}\left( \unit{J}\right) $ & $\mathbf{\Delta E}_{int}\left( \unit{J}%
\right) $ & $\mathbf{w}_{eng}\left( \unit{J}\right) $ & $\mathbf{Q}_{h}$ & $%
\mathbf{Q}_{c}$ \\ \hline
$\mathbf{IA}$ & $\left( 109.\,\allowbreak 38\unit{J}\right) $ & $\sim 0$ & $%
\left( 109.\,\allowbreak 38\unit{J}\right) $ & $\sim 0$ & $\left(
-109.\,\allowbreak 38\unit{J}\right) $ & $0$ \\ \hline
\end{tabular}%
\]%
But there is a difference. The new gas has useful fuel. and that means we
are ready to complete the cycle again by compressing the gas, igniting it,
and receiving the power stroke and exhaust. Note that Processes $AI$ and $IA$
have a net zero effect thermodynamically, but are very important to making
your car or lawn mower go!

\end{document}
